\scanpage{310}
\mspage{(9)}
qu'on donne la préférence \unsure{aux}\deleted{\unsure{l'un|chaines}} cochaines ou aux chaines. (Je n'ai pas eu un besoin urgent d'une telle convention de notation\ldots)

Dans le cas d'une dualité
\[
 v:A\overset{\approx}{\rightleftarrows}B^o
\]
on étend les conventions de montée et descente des indices aux objets à plusieurs indices, ainsi
\[
 (K_*{}^{.})^\vee=K^*{}_.\ \deleted{\unsure{\}}}.
\]
Dans la suite d'objets mixtes $(*)$ de la page précédente, on a écrit l'un en dessous de l'autre des objets qui se correspondent par dualité.

\uline{Remarque} Il est essentiel pour notre système de conventions de définir (disons) des bicomplexes de cochaines (objets de $A^{..}$) avec deux op \unsure{diff.|de diff.} $d^{\deleted{\unsure{(1)}}\, .}_{(1)}$ et $d^{.}_{(2)}$ qui \uline{commutent} -- on n'utilisera donc pas les conventions de Cartan-Eilenberg. Il faudra donc un signe à la diff. totale.

Du point de vue de la prop. 1 et du yoga \unsure{correspondant} ``description d'invariants associés à des $X_*\in\mathrm{Ob}\,\mathrm{Ss}$'', on peut dire que le \unsure{topoïs|topos|topo|théorie} de Dold-Puppe
% archive sheet number: 311
